% =====================================================================
%  PAGES LIMINAIRES : avertissement, registre des versions, préambule
% =====================================================================
\renewcommand{\lspdcourant}{Pages liminaires}

{\lspddisp\fontsize{22}{26}\selectfont\color{lspdnavy}Avertissement au lecteur}\\[4pt]
\lspdrulewide\vspace{3pt}
\noindent\begin{tikzpicture}\fill[lspdgold] (0,0) rectangle (2.6cm,.9pt);\end{tikzpicture}
\vspace{13pt}

\justifying
Le présent recueil constitue le \textbf{règlement intérieur} du Los Santos Police Department. Il rassemble en un document unique les règles d'organisation du service, les obligations déontologiques des personnels, les procédures opérationnelles de référence et le régime disciplinaire applicable.

\vspace{7pt}
Il s'impose à \textbf{l'ensemble des personnels} du département (agents assermentés comme personnels administratifs), sans considération de grade, d'ancienneté, d'affectation ou de division. Nul ne peut se prévaloir de son ignorance pour justifier un manquement.

\vspace{7pt}
Ce règlement \textbf{prime sur toute pratique orale, coutume de poste ou consigne informelle} qui lui serait contraire. Les notes de service émises par l'État-Major le complètent sans jamais le contredire ; en cas de contradiction apparente, la disposition du présent règlement prévaut jusqu'à modification formelle de celui-ci.

\vspace{7pt}
Il est susceptible d'être modifié à tout moment. Il appartient à chaque agent de s'assurer qu'il détient la version en vigueur et d'en connaître les évolutions.

\begin{note}[Comment lire ce règlement]
Chaque disposition porte un numéro d'alinéa de la forme \code{[12.04]} : le premier nombre désigne l'article, le second l'alinéa au sein de cet article. C'est sous cette référence qu'une disposition est citée dans un rapport, une demande de sanction ou un recours. Les encadrés bleus apportent une précision, les encadrés rouges signalent une interdiction ou une faute lourde, les encadrés verts rappellent une bonne pratique.
\end{note}

\lspdornament

\sstitre{Registre des versions}

\noindent
\begin{tabular}{P{24mm} P{26mm} P{34mm} P{60mm}}
\headrow
\thead{Version} & \thead{Date} & \thead{Autorité} & \thead{Objet} \\
\tcell{\textbf{1.0}} & \tcell{\textbf{\DATEVIGUEUR}} & \tcell{État-Major} & \tcell{Rédaction initiale et entrée en vigueur du règlement intérieur du département.} \\
\bottomrule
\end{tabular}

\vspace{4pt}
{\lspdsans\fontsize{8}{10}\selectfont\color{lspdgrey}
Toute modification ultérieure fait l'objet d'une nouvelle ligne au présent registre et d'une note de service diffusée à l'ensemble des personnels.}

\clearpage

% ---------------------------------------------------------------------
{\lspddisp\fontsize{22}{26}\selectfont\color{lspdnavy}Préambule}\\[4pt]
\lspdrulewide\vspace{3pt}
\noindent\begin{tikzpicture}\fill[lspdgold] (0,0) rectangle (2.6cm,.9pt);\end{tikzpicture}
\vspace{13pt}

\justifying
En qualité de service public de police, le Los Santos Police Department a pour mission la prévention des crimes, des délits et de l'ensemble des infractions aux lois en vigueur dans l'État de San Andreas, la protection des personnes et des biens, ainsi que le maintien de la paix publique sur le territoire du comté de Los Santos.

\vspace{7pt}
Cette mission ne se réduit pas à l'interpellation. Elle repose d'abord sur la \textbf{préservation de la vie humaine}, qui prime en toute circonstance sur l'objectif d'arrestation, sur la préservation des biens et sur toute considération tactique. Un agent qui met fin à une intervention sans interpellation mais sans qu'aucune vie n'ait été mise en péril a rempli sa mission.

\vspace{7pt}
L'autorité que la loi confie à l'agent de police n'est pas un privilège personnel : c'est un dépôt consenti par la population. Elle s'exerce sous le regard des citoyens, sous le contrôle de la hiérarchie et sous celui de l'autorité judiciaire. Chaque agent en répond individuellement.

\vspace{7pt}
C'est pour garantir cette exigence que le présent règlement a été établi. Il fixe ce que le département attend de chacun de ses membres, et ce que chacun de ses membres est en droit d'attendre du département.

\begin{solennel}[Devise du département]
\centering\normalsize\itshape
{\lspddisp\fontsize{13}{16}\selectfont\color{lspdnavy}
\upshape\addfontfeature{LetterSpace=12}\MakeUppercase{To Protect and to Serve}}\\[4pt]
\upshape\color{lspdgrey}\small Protéger et servir
\end{solennel}

\sstitre{Signalement des manquements}

Si vous constatez qu'un personnel du département contrevient directement au présent règlement, il vous appartient de le signaler à votre chaîne hiérarchique ou, lorsque la gravité des faits le justifie, à la Division des Affaires Internes. Le silence gardé sur une faute dont on a été témoin constitue lui-même un manquement.

\begin{rappel}[Ce que le département garantit à ses agents]
\begin{itemize}[leftmargin=6mm,topsep=2pt]
  \item une formation initiale et continue, et un encadrement identifié ;
  \item un équipement conforme et opérationnel, fourni par le département ;
  \item la protection du service en cas de menace, de harcèlement ou d'agression ;
  \item un accompagnement confidentiel en cas de difficulté psychologique ;
  \item le respect des droits de la défense en cas de procédure disciplinaire.
\end{itemize}
\end{rappel}

\clearpage
